%==============================================================================
% Appendix -- Extended in-silico metrics
%==============================================================================
\chapter{Extended in-silico metrics}
\label{app:extended}

This appendix collects the metric definitions, implementation details and full
metric distributions deferred from Chapters~\ref{ch:methods}
and~\ref{ch:results}, together with the octahedral scaling sweep of
Section~\ref{app:O24_scaling_and_metrics}. Each box plot shows the median,
interquartile range and whiskers over the design set, with the
lumazine-synthase reference (PDB~1NQX, dotted) shown where applicable.

\section{Metric definitions}
\label{app:metric-defs}
\label{app:designability}% alias: Figure~\ref{fig:octahedral} cites the metric
                         % definitions under this name.

All quantities are computed directly from the generated all-atom coordinates of
the symmetrised assembly; no external structure-prediction oracle is used. The
metrics fall into two families.

\paragraph{Backbone sanity.}
\emph{Chain breaks} are consecutive C$\alpha$--C$\alpha$ pairs whose bond length
departs from the canonical \SI{3.8}{\angstrom} by more than
$\tau_\text{cb}=\SI{0.75}{\angstrom}$; pairs spanning an intended chain
transition are masked out. The related \emph{maximum C$\alpha$ deviation} is the
worst such departure in a design, reported as a continuous summary rather than a
count. \emph{Backbone clashes} are heavy-atom pairs drawn from residues at least
two apart in sequence and closer than $\tau_\text{clash}=\SI{1.5}{\angstrom}$,
restricted to the backbone atoms $\{\text{N},\text{C}\alpha,\text{C}\}$, which
have no rotameric freedom with which to relieve an overlap; the corresponding
side-chain variant is much noisier and is reported only as a lower bound on
clash density. The \emph{non-loop fraction} is the proportion of residues
assigned to a helix or a strand by the P-SEA annotator
\parencite{labesse_p-sea_1997}, as implemented in Biotite's
\texttt{annotate\_sse}. As secondary filters we also report the radius of
gyration, the separate helix and sheet fractions, the number of
secondary-structure elements from the same annotation, and the per-residue
amino-acid composition, in particular the alanine and glycine content, where a
biased composition often signals a pathological design.

\paragraph{Symmetric-interface sanity.}
These use the full symmetrised complex, excluding any atom belonging to a fixed
or unsymmetrised motif, and all distances are
C$\alpha$--C$\alpha$. Three thresholds apply: a hard inter-subunit clash distance
of \SI{3.5}{\angstrom}, a contact band of \SIrange{4.0}{10.0}{\angstrom}, and a
proximity cutoff of \SI{15.0}{\angstrom} beyond which two subunits count as
non-interacting. \emph{ASU clashes} count inter-subunit pairs below the clash
distance, normalised by the number of subunits so that the figure is comparable
across point groups of different order. \emph{Mean contacts per interface}
counts, for every pair of subunits within the proximity cutoff, the pairs falling
inside the contact band, averaged over those interfaces. \emph{Minimum
inter-chain distance} is the smallest such distance anywhere in the complex:
values below roughly \SI{3.5}{\angstrom} indicate steric overlap, while values
well above \SI{6}{\angstrom} indicate that the subunits never meet, which for a
closed particle is itself a failure. The \emph{total interface contacts} is the
unnormalised sum of the contact counts over all interfaces, while the
\emph{number of interfaces with contacts} and the \emph{minimum contacts per
interface} together flag assemblies in which one interface is effectively absent
even though the mean is healthy. Finally, the \emph{minimum intra-ASU distance}
is the smallest C$\alpha$--C$\alpha$ distance \emph{within} a single subunit,
which catches intra-subunit overlap that the inter-subunit clash count cannot
see.

\section{Implementation notes}
\label{app:implementation}

\paragraph{Per-step symmetrisation.}
At each denoising step within the symmetrised portion of the trajectory:
(i)~Gaussian noise is added to the current coordinates and the noised structure is
itself symmetrised, so the network receives an exactly symmetric input
$\mathbf{X}^{(t)}$; (ii)~the network predicts clean coordinates
$\hat{\mathbf{X}}_0$ for all chains, which need not be exactly symmetric;
(iii)~the prediction is centred, the ASU slice is extracted and centred, and every
subunit (including the ASU itself) is rebuilt as
$R_g\,\hat{\mathbf{X}}_{0,\text{ASU}} + \lambda\,\mathbf{t}_g$, where
$\{(R_g,\mathbf{t}_g)\}_{g=1}^{G}$ are the rigid-body transforms of the chosen
point group and $\lambda$ is a scale factor that adapts the inter-subunit distance
to the current prediction; in the second half of the trajectory, a short
Lennard-Jones-based rigid-body optimisation of the ASU against its symmetry copies
is applied before this step to reduce inter-subunit clashes; (iv)~this symmetrised
prediction is used in the EDM update \parencite{karras_elucidating_2022}. For the
final $10\%$ of steps, no symmetry is enforced, allowing the model to relax any
residual strain at the inter-subunit interfaces before output.

\paragraph{Chunk distribution.}
Device $p$ receives $I_p$ rows as in Equation~\eqref{eq:row-stripe}, starting at
$p\lfloor I/P\rfloor + \min(p,\, I \bmod P)$. Uneven chunks are handled by padding
each local tensor to the per-rank maximum before the collective and trimming the
concatenated result back to $I$.

\paragraph{Communication volume.}
The dominant cost of the 1D scheme is the eighteen \texttt{ALLGATHER} calls on the
token track issued by the diffusion transformer per recycling iteration, with two
recycling iterations per denoising step. Each rank's message is sized by the token
count and channel width and is therefore independent of $P$, so the term that
grows with the device count is collective latency rather than payload.

\paragraph{Determinism.}
The two schemes reach determinism by different routes. In the 1D scheme the
atom-to-token pooling uses an index reduction whose floating-point accumulation
order depends on device scheduling, so ranks disagree slightly; because those
pooled features become the keys and values of the following cross-attention, the
discrepancy is amplified and breaks the requirement that all ranks hold identical
keys and values. A single broadcast of the pooled features immediately after the
reduction restores the invariant. The 2D scheme implements the same pooling as a
distributed scatter-reduce collective, which is deterministic by construction and
needs no broadcast. The two schemes also keep the sampling loop itself consistent
by different routes: the 1D scheme broadcasts the sampled noise and the denoised
prediction from rank~$0$ at each step boundary, whereas the 2D scheme draws the
full-shape noise tensor from per-rank generators held in lockstep and shards it
afterwards, so that no broadcast is required. Either way, a given seed produces
the same trajectory under either scheme.

\paragraph{Distributed nearest neighbours.}
The sparse atom attention needs the $k$ nearest neighbours of every query atom,
which a naive implementation obtains from an $L\times L$ distance matrix, the very
object that cannot be afforded at these scales. Instead, each device computes
distances between its local atom rows and those resident on its column partner,
in chunks, so that even the per-quadrant block is never held in full; a ring
reduction then merges partial top-$k$ candidates around the column axis over
$\sqrt{P}$ steps. After the final step each device holds the global top-$k$
indices for its own query atoms, and no $L\times L$ tensor is materialised on any
device at any point.

\clearpage
\section{Icosahedral designs}
\label{app:ico}

The comparison in Figure~\ref{fig:icosahedral}b--e covers four panels: chain
breaks, backbone clashes, non-loop fraction and maximum C$\alpha$ deviation.
Figure~\ref{fig:ico-perchain} shows the full eight-panel version, adding helix
fraction, sheet fraction, glycine content and the average number of
secondary-structure elements per chain.

\paragraph{Caveat on problem difficulty.}
The two design problems are not of equal difficulty, and the comparison should be
read with this asymmetry in mind. Each icosahedral chain is denoised inside a
$12{,}600$-residue joint context, where its trajectory must remain
self-consistent with the simultaneously denoised trajectories of $59$ symmetry
mates and with all the inter-chain interfaces they form, whereas the monomer
problem is a single $210$-residue chain comfortably inside RFD3's native
$384$-token training crop. The icosahedral problem is therefore strictly harder
per chain.

\paragraph{Helix and sheet fraction.}
The most prominent qualitative gap in Figure~\ref{fig:ico-perchain} appears in
secondary-structure composition. Vanilla RFD3 monomers reproduce a helix/sheet
balance close to that of the 1NQX reference: the helix fraction concentrates
around a median of $\approx 0.47$ with a tight bulk against the 1NQX value of
$0.47$, and the sheet fraction sits at a median of $\approx 0.18$ with a
comparable spread, against $0.20$ for 1NQX. The Design-CP icosahedral designs, in
contrast, are strongly $\beta$-biased: the helix fraction collapses near zero on
the bulk of chains, with a few upper outliers, and the sheet fraction shifts
upwards to a median of $\approx 0.5$ with a noticeably broader distribution.

\paragraph{Number of secondary-structure elements.}
Consistent with the helix/sheet shift, the average number of secondary-structure
elements per chain rises from $\approx 13$ in the monomers to $\approx 16$ in the
icosahedral designs, with the icosahedral distribution carrying a heavier upper
tail; both populations include outliers, but only the icosahedral set reaches the
upper twenties. At fixed chain length a higher element count mechanically implies
shorter average elements, which is again consistent with the icosahedral chains
preferring multiple short $\beta$-strands over a small number of long helices. We
treat this as supportive evidence for the helix/sheet shift rather than an
independent observation, and note that the metric is sensitive to the P-SEA
assignment thresholds \parencite{labesse_p-sea_1997}.

\paragraph{Amino-acid composition.}
On amino-acid composition the two populations are closer to each other than on
secondary structure. Alanine content is similar across both sets, with broadly
overlapping distributions and medians of the same order, and is, as is typical
for RFD3 \parencite{butcher_novo_2025}, slightly inflated relative to 1NQX in
both cases; we do not read a meaningful difference between Design-CP ASUs and
vanilla monomers on this axis. Glycine content, in contrast, is appreciably
broader in the icosahedral designs ($\approx 0.10$--$0.25$, with sporadic high
outliers) than in the monomers ($\approx 0.05$--$0.10$, very tightly
concentrated), although the medians remain of the same order and within the
typical range for natural proteins. The widened glycine spread is qualitatively
consistent at the population level with the broader maximum
C$\alpha$--C$\alpha$ deviation of Figure~\ref{fig:icosahedral}e, chains under
greater inter-subunit strain relying more on backbone-flexible residues, but we
flag this as a tentative association rather than a causal claim, since we have
not tested whether the same chains drive both effects.

\begin{figure}[h]
  \centering
  \includegraphics[width=\textwidth]{RFD3-vs-DesignCP.png}
  \caption{\textbf{Per-chain comparison of Design-CP icosahedral designs against
    vanilla RFD3 monomers (full eight panels).} Backbone-sanity and composition
    metrics computed chain by chain on $n{=}40$ Design-CP icosahedral assemblies
    ($60$ chains $\times\,210$ residues) and $n{=}40$ vanilla single-GPU RFD3
    monomers of length $210$, with the lumazine-synthase reference (PDB~1NQX)
    dotted in each panel. The first four panels reproduce
    Figure~\ref{fig:icosahedral}\textbf{b--e}; the additional four (helix
    fraction, sheet fraction, glycine content, average number of
    secondary-structure elements) show the secondary-structure shift discussed
    above.}
  \label{fig:ico-perchain}
\end{figure}

\begin{figure}[h]
  \centering
  \includegraphics[width=\textwidth]{Icosahedral_designability_metric_ASU.png}
  \caption{Full designability (backbone-sanity) distributions for the icosahedral
    I60 designs ($n{=}40$), including radius of gyration.}
  \label{fig:ico-des}
\end{figure}

\begin{figure}[h]
  \centering
  \includegraphics[width=\textwidth]{Icosahedral_symmetry_metrics_extra.png}
  \caption{Additional symmetry-aware interface metrics for the icosahedral I60
    designs ($n{=}40$), extending Figure~\ref{fig:icosahedral}\textbf{g--l}.}
  \label{fig:ico-symm}
\end{figure}

\clearpage
\section{Octahedral designs on workstation-grade GPUs}
\label{app:O24_scaling_and_metrics}
\label{app:oct}

This section collects the material deferred from
Section~\ref{sec:results-democratization}: the scaling sweep for octahedral
targets on workstation-grade hardware, and the metric distributions not shown in
the main-text figure.

\subsection*{Maximum assembly size on RTX~A4000 GPUs}

The sweep follows the protocol of Section~\ref{sec:results-scaling}: at a fixed
device count we increase the ASU length until inference runs out of memory, and
record the last feasible length. The only difference is the hardware. Each
RTX~A4000 provides $16$\,GB of device memory against the
$95$\,GB of a GH200, roughly a sixfold reduction per device, so even
the aggregate memory of the whole A4000 cluster is below that of a single GH200.
Because octahedral symmetry has order $24$, a given ASU length here corresponds
to an assembly of $24$ chains rather than the $60$ of the icosahedral case, and
the token and atom counts entering the quadratic pair tracks scale accordingly.

On $16$ A4000s the 2D scheme reached a ceiling of $178$ residues per chain, that
is $4{,}272$ residues over the full assembly, and this is the configuration used
for the $n{=}12$ designs analysed below. Two points follow. First, the ceiling is
set by aggregate rather than per-device memory, so the sixfold per-card deficit
is largely recovered by the device count, and the per-chain sizes reached remain
comparable to those of existing \emph{de novo} octahedral nanoparticles
\parencite{king_computational_2012}. Second, the 2D scheme requires a
perfect-square device count, so on a cluster of this size the usable
configurations are $P\in\{1,4,9,16\}$ and $16$ is the largest available; the 1D
scheme carries no such restriction and is the more practical choice at
intermediate device counts on commodity hardware.

\subsection*{Full metric distributions}

Figures~\ref{fig:oct-des} and~\ref{fig:oct-symm} give the complete distributions
over the $n{=}12$ octahedral designs for both metric families of
Section~\ref{app:metric-defs}. The headline panels of
Figure~\ref{fig:octahedral}b--g are reproduced there; we comment below only on
the additional panels.

\paragraph{Additional backbone-sanity panels.}
The radius of gyration is tightly concentrated in the
\SIrange{56.8}{57.4}{\angstrom} range, indicating a consistent per-chain envelope
across the population. Helix and sheet fractions are both broad, with no
dominant secondary-structure type within the population: the helix fraction
spans $\approx 0$--$0.9$ with a median around $\approx 0.55$, and the sheet
fraction spans $\approx 0$--$0.85$ with a median around $\approx 0.45$. This
contrasts qualitatively with the icosahedral designs of
Section~\ref{app:ico}, where the helix fraction collapses near zero. We interpret
this as plausibly reflecting the lower oligomeric state ($24$ against $60$
subunits), which gives the symmetrisation step less reason to favour extended
$\beta$-sheets over helical packing, but caution that the sample size
($n{=}12$) is small. Alanine content has a median of $\approx 0.25$ and is
comparatively broad ($\approx 0.05$--$0.6$), while glycine content is tightly
clustered around $\approx 0.05$--$0.10$, in line with the icosahedral population.

\paragraph{Additional symmetry-interface panels.}
The hard inter-subunit clash counts at both the ASU and complex levels are zero
across the entire population, indicating that no design realises sterically
forbidden inter-subunit contacts. The smallest intra-ASU C$\alpha$--C$\alpha$
distance is centred at $\approx\SI{3.6}{\angstrom}$, matching the canonical
C$\alpha$--C$\alpha$ spacing, with two low outliers at $\approx\SI{1.2}{\angstrom}$
and $\approx\SI{2.4}{\angstrom}$ that flag designs with intra-ASU geometric
strain. The number of interfaces showing contacts has a median of $\approx 55$,
and the total number of interface contacts spans $\approx 200$--$3{,}700$ with a
median of $\approx 1{,}700$. The minimum number of contacts per interface has a
median near $1$, with two outliers at $\approx 25$ and $\approx 30$: some designs
do contain a weak interface whose contact count drags the per-design minimum
towards zero, which is a useful filter signal for downstream design campaigns.

Across both families, the additional panels are consistent with the headline
result of Section~\ref{sec:results-democratization}: octahedral designs sampled
on commodity GPUs satisfy the same hard-failure criteria as the icosahedral
baseline, the main qualitative difference being a more permissive
secondary-structure distribution.

\begin{figure}[h]
  \centering
  \includegraphics[width=\textwidth]{Octahedral_designability_metrics.png}
  \caption{Full designability (backbone-sanity) distributions for the octahedral
    O24 designs ($n{=}12$).}
  \label{fig:oct-des}
\end{figure}

\begin{figure}[h]
  \centering
  \includegraphics[width=\textwidth]{Octahedral_symmetry_metrics.png}
  \caption{Full symmetry-aware interface metric distributions for the octahedral
    O24 designs ($n{=}12$).}
  \label{fig:oct-symm}
\end{figure}

\clearpage
\section{Effect of removing the symmetry constraint}
\label{app:monstromers}

\begin{figure}[h]
  \centering
  \includegraphics[width=\textwidth]{Monstromers.png}
  \caption{Six Design-CP samples at the icosahedral system size ($60$ chains
    $\times\,210$ residues $=12{,}600$ residues) generated \emph{without} a
    point-group symmetry constraint. Compared with the well-formed icosahedral
    particles of Figure~\ref{fig:icosahedral}, the unconstrained assemblies are
    malformed and non-globular, illustrating why the symmetry prior is essential in
    this beyond-crop regime.}
  \label{fig:monstromers}
\end{figure}
